\section{Introdução}

Todo sistema é desenvolvido com um propósito. A forma de saber se o sistema está atendendo ao seu propósito é estabelecendo requisitos que impõem restrições e direcionam o processo de desenvolvimento para a direção correta. Os requisitos em si são uma grande área de pesquisa, mas podem ser principalmente divididos em 2 grupos: funcionais e não funcionais.

Requisitos funcionais são restrições que especificam o que o sistema deve ser capaz de fazer. Um requisito funcional para uma rede social pode ser a capacidade de postar fotos. Por outro lado, requisitos não funcionais são restrições que definem como algo deve ser feito. Um requisito não funcional para a rede social mencionada pode ser o armazenamento seguro das fotos dos usuários, impedindo acesso não autorizado a elas. Em resumo, requisitos funcionais definem funcionalidades, enquanto os não funcionais estabelecem expectativas de qualidade.

A arquitetura de software é definida como uma coleção de decisões de design que afetam a estrutura, comportamento e qualidade geral de um sistema de software, servindo como base para decisões subsequentes. É o campo de pesquisa que lança luz sobre requisitos não funcionais. Ao longo dos anos, muitas dessas decisões de design de qualidade foram agrupadas no que são chamados de padrões arquiteturais ou simplesmente padrões.

Assim, padrões podem ser definidos como conceitos reutilizáveis que podem ser aplicados para resolver questões similares em sistemas em contextos similares, considerando necessidades não funcionais. Isso significa que um padrão pode ser aplicado a qualquer sistema, independentemente de suas funcionalidades, desde que tenham requisitos de qualidade similares.

Se funcionalidades fossem a única parte importante, não haveria necessidade de ter arquitetura de software. Mas ao mesmo tempo, isso não é o que acontece na realidade. Requisitos de qualidade acabam não sendo tão relevantes quanto deveriam ao escolher padrões arquiteturais para um sistema na indústria, e isso traz consequências de longo prazo para confiabilidade, manutenibilidade, testabilidade, entre outras complicações.

O propósito deste trabalho é experimentar com um padrão arquitetural específico chamado Modelo de Atores. A ideia é usar este padrão em um projeto que está bem distante dos requisitos de qualidade que ele atende e entender quais são as consequências.
